% !PRINTED 373
\weberpage{373}
%<<CONT
% Source-faithful Cayley configuration signs: three/four concurrent strokes and a seven-rayed star.
\providecommand{\weberthreeup}{\mathord{\tikz[baseline=-0.2ex,x=0.23em,y=0.23em,line width=0.45pt]{\draw (0,0)--(-1,1) (0,0)--(0,1.08) (0,0)--(1,1);}}}
\providecommand{\weberfourup}{\mathord{\tikz[baseline=-0.2ex,x=0.23em,y=0.23em,line width=0.45pt]{\draw (0,0)--(-1.05,1) (0,0)--(-0.35,1.08) (0,0)--(0.35,1.08) (0,0)--(1.05,1);}}}
\providecommand{\webersevenstar}{\mathord{\tikz[baseline=-0.35ex,x=0.21em,y=0.21em,line width=0.45pt]{\draw (0,0)--(90:1) (0,0)--(38.57:1) (0,0)--(-12.86:1) (0,0)--(-64.29:1) (0,0)--(-115.71:1) (0,0)--(-167.14:1) (0,0)--(141.43:1);}}}
\gesperrt{The quadruples of double tangents with the
signs}
\[
\mid\mid\mid\mid,\ \square
\]
\gesperrt{are syzygetic, and those with the signs}
\[
\weberfourup,\ \triangle\mid,\ \weberthreeup\mid,\ \vee\vee
\]
\gesperrt{azygetic.}

All remaining quadruples are neither syzygetic nor azygetic.

After this it is easy to form the signs for all the complete
seven-systems.

Such a sign must consist of seven strokes, which cannot have
more than eight vertices and which form a figure
from which neither of the two figures $\mid\mid\mid$, $\sqcap$ can be
detached. From this it follows first that this figure can consist of not more than
two separate parts, because otherwise the figure $\mid\mid\mid$
would be contained in it, and that no part can have more than one centre
from which several strokes issue, except when
this part is the triangle $\triangle$, because otherwise the figure $\sqcap$ would
occur.

Now if the figure is in one part, it must be a seven-rayed
star $\webersevenstar$, and since one can choose any of the eight figures as
middle point, these are eight possibilities, of
which one is the assumption considered above:
\[
[1\,8]\,[2\,8]\,[3\,8]\,[4\,8]\,[5\,8]\,[6\,8]\,[7\,8].
\]
But if the figure is in two parts, it is first to be excluded
that the one part consists of one or of two strokes or of
a three-rayed star; for in these cases the other part
would have to be a star with six, five or four rays.
But for that not enough of the eight figures remain.
It therefore only remains that the one part is a
triangle, the other a four-rayed star, $\triangle\,\weberfourup$, and this
assumption is also in fact always admissible. A representative
of such a system is:
\[
[1\,2]\,[1\,3]\,[2\,3]\,[4\,5]\,[4\,6]\,[4\,7]\,[4\,8].
\]
The triangle can be chosen in $8.7.6 : 2.3 = 56$ different
ways, and for every triangle there are still five possibilities
of taking the four-rayed star, since one can make any of the
remaining five points the centre. The number
%CONT>>
